\section{Desarrollo e implementación del backend basado en Blosc}
\label{sec:development}

Esta sección presenta el desarrollo del mecanismo de almacenamiento comprimido basado en
Blosc incorporado en TMFQSfullstate. La descripción se organiza alrededor de la
disposición en bloques del vector de estado, la caché LRU de bloques descomprimidos y la
estrategia de acceso local utilizada durante la simulación.

\begin{figure}[H]
    \centering
    \fbox{\parbox{0.82\linewidth}{\centering
        Espacio reservado para un diagrama de módulos del desarrollo realizado, mostrando la relación entre el registro cuántico, el backend Blosc,
        la disposición en bloques, la caché y el flujo de compresión/descompresión.}}
    \caption{Arquitectura del desarrollo implementado para el backend basado en Blosc}
    \label{fig:development-architecture}
\end{figure}

\subsection{Contexto de implementación}

La integración se realizó sobre TMFQSfullstate, manteniendo inalterada la interfaz lógica del simulador y concentrando los cambios en la capa de
almacenamiento del vector de estado. Esta decisión permite que las operaciones de más alto nivel sigan trabajando con amplitudes complejas, compuertas
y mediciones, mientras que la representación física del estado puede cambiar de densa a comprimida sin afectar la semántica del simulador.

En ese contexto, el backend desarrollado se apoya en tres decisiones concretas. La primera es representar el estado mediante bloques de amplitudes.
La segunda es emplear Blosc2 como biblioteca de compresión sin pérdida para cada bloque. La tercera es introducir una caché LRU de bloques
descomprimidos para amortiguar accesos repetidos durante la aplicación de compuertas.

\subsection{Backend de almacenamiento comprimido basado en Blosc}

La implementación realizada adopta Blosc2 como realización concreta de la estrategia abstracta descrita antes. El estado cuántico completo se divide
en bloques de tamaño configurable, y cada bloque se almacena de forma comprimida. Esta organización permite que la descompresión ocurra sobre una
unidad local y no sobre el vector completo.

Una ventaja importante de esta elección es que Blosc2 proporciona soporte directo para compresión rápida por bloques, así como filtros previos que
mejoran la compresibilidad de arreglos numéricos. En el contexto del simulador, esto se traduce en una mejor adecuación al patrón ``descomprimir
solo lo necesario, modificar localmente y recomprimir cuando corresponda''.

Además, el backend expone parámetros de configuración relevantes para el experimento: número de estados por bloque, nivel de compresión, número de
hilos y cantidad de entradas de caché. Estos parámetros permiten ajustar el compromiso entre uso de memoria y costo temporal, y por ello deben ser
reportados explícitamente cuando se presenten resultados experimentales.

\subsection{Disposición en bloques y mapeo del vector de estado}

La primera parte del desarrollo consistió en introducir una disposición explícita del vector de estado en bloques contiguos. Cada bloque agrupa una
cantidad fija de estados base y almacena sus componentes real e imaginaria en forma intercalada. Con ello, cualquier estado base puede mapearse a un
índice de bloque y a un desplazamiento local dentro de ese bloque.

Esta capa de mapeo es necesaria para que el backend identifique rápidamente qué bloque debe adquirirse ante una operación de lectura, escritura o
aplicación de compuerta. También permite tratar el último bloque de manera especial cuando el número total de estados no llena exactamente una unidad
completa de almacenamiento.

\begin{figure}[H]
    \centering
    \fbox{\parbox{0.82\linewidth}{\centering
        Espacio reservado para una figura que muestre el mapeo entre índices globales del vector de estado, bloques físicos y desplazamientos locales
        dentro de cada bloque.}}
    \caption{Mapeo entre estados base y bloques físicos en el backend comprimido}
    \label{fig:development-layout}
\end{figure}

\subsection{Caché LRU de bloques descomprimidos}

El segundo componente clave del desarrollo fue la incorporación de una caché de bloques descomprimidos. Su objetivo es mantener en memoria de trabajo
los bloques usados recientemente, de modo que múltiples operaciones sucesivas puedan reutilizarlos sin repetir el proceso de descompresión.

La política de reemplazo adoptada es LRU. Cada acceso actualiza la marca de uso reciente del bloque, y cuando la caché se llena se selecciona para
expulsión el bloque con el acceso más antiguo. Si el bloque expulsado fue modificado, se marca como sucio y se recomprime antes de salir de la caché;
si no hubo modificaciones, puede descartarse sin costo adicional de escritura.

Desde el punto de vista de la simulación, esta decisión es importante porque muchas secuencias de compuertas reutilizan un subconjunto reducido del
estado. En tales casos, una caché pequeña pero bien administrada puede reducir de forma significativa la frecuencia de descompresiones y recompresiones.

\subsection{Estrategia de acceso y actualización local}

La implementación concreta sigue el siguiente esquema operativo. Cuando una operación necesita una amplitud o un par de amplitudes, primero determina
los bloques implicados. Después consulta la caché. Si el bloque ya está disponible, reutiliza el buffer descomprimido. Si no está disponible, lo
recupera desde almacenamiento comprimido, lo inserta en la caché y continúa con la operación.

Una vez modificado un bloque, este no se recomprime inmediatamente por defecto. En lugar de ello, se mantiene en la caché marcado como sucio hasta
que sea necesario liberarlo o hasta que la operación actual requiera un vaciado explícito. Esta política evita recompresiones prematuras cuando varias
compuertas consecutivas afectan el mismo conjunto de bloques.

Para compuertas cuyas amplitudes involucradas residen en un mismo bloque, el procedimiento es completamente local. En cambio, para operaciones que
afectan amplitudes distribuidas en bloques distintos, el backend debe adquirir ambos bloques y gestionar su permanencia en la caché. Este caso hace
evidente la necesidad de al menos dos ranuras activas y explica por qué la caché no puede reducirse a una sola entrada.

\begin{figure}[H]
    \centering
    \fbox{\parbox{0.82\linewidth}{\centering
        Espacio reservado para una secuencia de pasos de acceso: localizar bloque, verificar caché, descomprimir, modificar, marcar como sucio y
        eventualmente recomprimir al expulsar.}}
    \caption{Secuencia operativa del backend basado en Blosc}
    \label{fig:development-workflow}
\end{figure}

El mecanismo descrito en este capítulo establece la base operativa para medir el efecto de la
compresión sobre circuitos representativos. La sección siguiente presenta la evaluación
experimental realizada sobre Grover y QFT, así como la comparación cuantitativa entre la
estrategia sin compresión, la estrategia sin pérdida basada en Blosc y una estrategia con
pérdida controlada basada en ZFP.
